% Part: computability
% Chapter: tm-computations
% Section: church-turing-thesis

\documentclass[../../../include/open-logic-section]{subfiles}

\begin{document}

\olfileid[hi]{tur}{mac}{ctt}
\olsection{चर्च--ट्यूरिंग अभिधारणा}

ट्यूरिंग मशीनों को प्रभावी प्रक्रिया की धारणा का सटीक प्रतिस्थापन माना
जाता है। ट्यूरिंग का विचार था कि प्रभावी प्रक्रिया और ट्यूरिंग मशीन की
दोनों धारणाएँ समझने वाले प्रत्येक व्यक्ति को यह सहज बोध होगा कि प्रभावी
प्रक्रिया द्वारा किया जा सकने वाला प्रत्येक कार्य ट्यूरिंग मशीन द्वारा
भी किया जा सकता है। इस दावे को इस तथ्य से समर्थन मिलता है कि प्रभावी
प्रक्रिया की धारणा के लिए प्रस्तावित अन्य सभी सटीक प्रतिस्थापन ट्यूरिंग
मशीन की धारणा के विस्तारतः समतुल्य निकलते हैं---अर्थात् वे ठीक उसी
फलन-समुच्चय को संगणित कर सकते हैं। इस दावे को \emph{चर्च--ट्यूरिंग
अभिधारणा} कहते हैं।

\begin{defn}[चर्च--ट्यूरिंग अभिधारणा]
\emph{चर्च--ट्यूरिंग अभिधारणा} कहती है कि प्रभावी प्रक्रिया द्वारा
संगणनीय प्रत्येक वस्तु ट्यूरिंग-संगणनीय है।
\end{defn}

चर्च--ट्यूरिंग अभिधारणा का आह्वान दो प्रकार से किया जाता है। उसका पहला
प्रकार का उपयोग आलस्य का बहाना है। मान लें कि हमारे पास किसी वस्तु को
संगणित करने की प्रभावी प्रक्रिया का वर्णन, मान लें ``छद्म-कूट'' में,
है। तब चर्च--ट्यूरिंग अभिधारणा का आह्वान करके इस दावे को उचित ठहरा सकते
हैं कि वही फलन किसी ट्यूरिंग मशीन द्वारा संगणित होता है, भले ही हमने
वास्तव में उस मशीन का निर्माण न किया हो।

चर्च--ट्यूरिंग अभिधारणा का दूसरा उपयोग दार्शनिक दृष्टि से अधिक रोचक है।
यह दिखाया जा सकता है कि ऐसे फलन हैं जिन्हें ट्यूरिंग मशीनें संगणित नहीं
कर सकतीं। इससे चर्च--ट्यूरिंग अभिधारणा का उपयोग करके निष्कर्ष निकाला जा
सकता है कि किसी भी प्रक्रिया द्वारा उन्हें प्रभावी रूप से संगणित नहीं
किया जा सकता। क्योंकि यदि ऐसी प्रक्रिया होती, तो चर्च--ट्यूरिंग
अभिधारणा से उसके लिए ट्यूरिंग मशीन का अस्तित्व निकलता। अतः यदि हम सिद्ध
कर सकें कि उसे संगणित करने वाली कोई ट्यूरिंग मशीन नहीं है, तो कोई
प्रभावी प्रक्रिया भी नहीं हो सकती। विशेष रूप से, चर्च--ट्यूरिंग
अभिधारणा का आह्वान यह दावा करने के लिए किया जाता है कि तथाकथित विराम
समस्या को न केवल ट्यूरिंग मशीनों द्वारा हल नहीं किया जा सकता, बल्कि उसे
किसी भी प्रकार प्रभावी रूप से हल नहीं किया जा सकता।

\end{document}
